\section{Declaración de uso de inteligencia artificial}
\label{sec:declaracion-ia}

\begin{xltabular}{\textwidth}{@{} >{\raggedright\arraybackslash}p{2.2cm} >{\raggedright\arraybackslash}p{2.4cm} >{\raggedright\arraybackslash}p{1.7cm} >{\raggedright\arraybackslash}p{3.6cm} >{\raggedright\arraybackslash}p{3.8cm} @{}}
  \caption{Herramientas de inteligencia artificial usadas como apoyo.}
  \label{tab:declaracion-ia} \\
  \toprule
  \textbf{Integrante} & \textbf{Herramienta} & \textbf{Versión} & \textbf{Para qué se usó} & \textbf{Cómo se verificó lo obtenido} \\
  \midrule
  \endfirsthead
  \toprule
  \textbf{Integrante} & \textbf{Herramienta} & \textbf{Versión} & \textbf{Para qué se usó} & \textbf{Cómo se verificó lo obtenido} \\
  \midrule
  \endhead
  \bottomrule
  \endfoot

  Steven Pozo & OpenAI Codex & GPT-5 &
  Organización y redacción de borradores del resumen de la entrevista,
  contexto, historias de usuario y trabajos relacionados; búsqueda y
  revisión de referencias; apoyo inicial en la bitácora. &
  Contrastó los borradores con los datos entregados por el equipo, mantuvo
  como pendientes los datos no confirmados, verificó los DOI contra las
  páginas de las revistas y revisará el contenido con el equipo y la
  entrevistada antes de usarlo como evidencia. \\
  \addlinespace

  Dariem Benites & Claude Code (Anthropic) & Claude Opus 5.5 &
  Usé Claude para que me ayude a corregir los borradores del ADR-001 y
  de las secciones 06, 07, 08, 09 y 10a con la entrevista, los RNF y los
  riesgos; resumir el artículo de Blinowski et al.; pasar tablas a
  xltabular y compilar; revisar la coherencia con los RNF y los riesgos, y
  aplicar pequeñas correcciones a las aclaraciones del \#20. &
  Comparé los diagramas con mi bosquejo, la entrevista y las correcciones
  de Mario; revisé cada cifra contra la entrevista y la sección 05;
  comprobé las referencias con su DOI; compilé sin errores y revisé tablas
  y figuras en el PDF; y leí cada sección hasta poder explicarla. \\
  \addlinespace

  Jilmar Cedeño & OpenAI Codex & GPT-6.2 ligero &
  Usé Codex para preparar y ajustar borradores de objetivos, requisitos no
  funcionales y riesgos, y para comparar las secciones en la revisión del
  issue \#20. &
  Contrasté los datos de los requisitos no funcionales con la entrevista y
  comprobé los hallazgos en las secciones 08 y 09, el riesgo R02 y el
  ADR-001. \\
  \addlinespace

  Mario Castro & Claude (Anthropic) & Claude Opus 5.5 &
  Usé Claude como apoyo. La elección de gRPC y el alcance del esqueleto ya
  los habíamos decidido en el equipo (issues \#10 y \#11). Escribí
  una primera versión del código, basada en la de PE-U1, y borradores del
  ADR-002, de las secciones 10b, 12 y 13 y del README, y Claude corrigió el formato de las tablas de aportes y de declaración de IA (xltabular), resolvió los conflictos al fusionar las ramas del equipo y compiló el documento varias veces. En esas compilaciones encontró varios errores: la versión de Python estaba mal puesta, había un desborde de medio punto y un proceso se había quedado colgado en el puerto 50051. &
  Seguí solo el README en un clon limpio, en Windows~11, el 21 de noviembre
  de 2026, y comprobé los cuatro pasos, incluido el de Solicitudes
  apagado. Contrasté el ADR-002 con el código de sockets y gRPC de PE-U1 y
  con las secciones 08 y 09, y dejé pendientes las cifras de latencia hasta
  tener la corrida oficial. Compilé el documento y revisé las secciones en
  el PDF. \\
  \addlinespace

\end{xltabular}
